\section{Problem Statement and Method}
\label{s:method}

\subsection{Problem Statement}

In this section, we formulate the task as a supervised binary classification problem.
Given a 3D CT image $\bm{I}\in\mathbb{R}^{d \times h \times w}$ of a patient,
our goal is to learn a classifier $\xi$ which can precisely predict the residual disease $y$,
with $y = \xi(\bm{I})$.
No residual disease is known as R0 and vice versa as non-R0 in medical terms.
In this paper, we denote $y=1$ for R0 and $y=0$ for non-R0.

\subsection{Overall Architecture}
The overall architecture of our proposed method is shown in Fig.~\ref{f:framework}.
Firstly, we use pooling to create multi-view slice features from the various dimensions of 3D CT images.
Then, we develop a~\fullname~(\shortname) to learn more valuable image features under different views with an attention mechanism.
Finally, we integrate the learned multi-view features and employ Med3D to help the final label prediction.
In the following part, we will introduce~\shortname~in detail.
\begin{figure*}
    \centering
    \includegraphics[width=\textwidth]{figure/framework.pdf}
    \caption{The architecture of MuVAL.}
     \label{f:framework}
    % \vspace{-4mm}
\end{figure*}

\subsection{Multi-View Attention Learning}
As shown in Fig.~\ref{f:MuVAL_sample}, 
the raw CT image has the natural characteristics of a 3D structure.
Observing from different views can obtain different features and thus enhance the representation of 3D CT images.
For a 3D CT image $\bm{I}\in\mathbb{R}^{d\times h \times w}$,
we consider the transverse view as the default view and augment it with the coronal and sagittal views to obtain multiple views.
%from depth dimension~(\textit{i.e.}, $d$).
%To obtain multiple views, we leverage the coronal view and sagittal view
%the height dimension (\textit{i.e.}, $h$) and width dimension (\textit{i.e.}, $w$) 
%as the addition to the transverse view.
In addition, attention mechanism has achieved excellent results on various computer vision tasks~\cite{ruan2021color,ruan2021dae}.
Along this line, we further employ attention mechanism to identify the most relevant slices in each view, assigning different weights to different slices. 
%\textit{i.e.}, different slices should have different attention weights. 
To accomplish this, we apply the SE module proposed by Hu et al.\cite{hu2018squeeze}, which is typically used to obtain channel attention weights in 2D CNN. 
%to obtain the channel attention weights in 2D CNN. 
 In our case, we treat the slices as channels, and the SE module is used to obtain the slice attention weights. 
Inspired by this approach, we propose a Multi-View Attention Learning (MuVAL) method to obtain slice attention weights from multiple views. 

To construct multi-view slice features, we first apply global average pooling to different dimensions of 3D CT images, which can be formulated as follows: 
\begin{equation}
\begin{aligned}
%\begin{array}{ll}
    \bm{f}_t(i) = \frac{1}{h \times w}\sum_{j=1}^{h}\sum_{k=1}^{w}\bm{I}(i,j,k),\\
    \bm{f}_c(j) = \frac{1}{d \times w}\sum_{i=1}^{d}\sum_{k=1}^{w}\bm{I}(i,j,k),\\
    \bm{f}_s(k) = \frac{1}{d \times h}\sum_{i=1}^{d}\sum_{j=1}^{h}\bm{I}(i,j,k),\\
%\end{array}
\end{aligned}
\end{equation}
where $\bm{f}_t(i)$ denotes the $i$-th element of $\bm{f}_t$, which is generated by shrinking the $i$-th slice of $\bm{I}$ in the transverse view.

The slice features $\bm{f}_t$, $\bm{f}_c$ and $\bm{f}_s$, can be seen as a condensation of the slice spatial information for each view.
To further obtain the slice attention weights, we employ attention modules consisting of a bottleneck with two fully-connected (FC) layers around the ReLU\cite{nair2010rectified} function with a sigmoid activation:
\begin{equation}
\begin{aligned}
%\begin{array}{ll}
    \bm{\alpha}_t  = \sigma(\bm{W}_{t2}\delta(\bm{W}_{t1}\bm{f}_t)),\\
    \bm{\alpha}_c  = \sigma(\bm{W}_{c2}\delta(\bm{W}_{c1}\bm{f}_c)),\\
    \bm{\alpha}_s  = \sigma(\bm{W}_{s2}\delta(\bm{W}_{s1}\bm{f}_s)),
%\end{array}
\end{aligned}
\end{equation}
where $\sigma$ refers to the sigmoid activation, $\delta$ refers to the ReLU function, $\bm{W}_{t1} \in \mathbb{R}^{\frac{d}{r} \times d}$ and $\bm{W}_{t2} \in \mathbb{R}^{d \times \frac{d}{r}}$, $r$ is the dimensionality-reduction ratio.

To enhance the representation of 3D CT images, we apply the average to fuse the multi-view slice attention: 
\begin{equation}
\overline{\bm{I}}(i,j,k)  = \frac{\bm{\alpha}_t(i) + \bm{\alpha}_c(j) + \bm{\alpha}_s(k)}{3}\bm{I}(i,j,k).
\end{equation}

\subsection{Residual Disease Prediction}
For the classification stage, we apply the Med3D~\cite{chen2019med3d} which aggregated several 3D medical image datasets and obtained a series of pre-trained ResNet-3D models. This approach has demonstrated impressive results in various medical challenges and datasets.\cite{li2021alzheimer}.

%ResNet的图，结构说一下,max pool的好处
Considering the computational complexity and available data scale, we employ a pre-trained model of Med3D as the backbone, with a 34-layer ResNet-3D structure. For each input $\bm{\overline{I}}$, the backbone outputs $n$-dimensional features:

\begin{equation}
\bm{f} = Backbone(\overline{\bm{I}}), \bm{f} \in \mathbb{R}^n.
\end{equation}

To ultimately generate the classification probabilities for non-R0 and R0, we apply a 2-output FC layer with a softmax activation:
\begin{equation}
\bm{p} = \sigma(\bm{W}\bm{f}),
\end{equation}
where $\sigma$ refers to the softmax activation, $\bm{W} \in \mathbb{R}^{2 \times n}$, $\bm{p}$ is the one-hot representation for the classification probabilities.

\subsection{Model Learning}
For model learning, 
we employ the cross-entropy as the loss function since it is a classification problem.
The loss function for the output of the last layer is shown as follows:
\begin{equation}
L = -\frac{1}{N}\sum_{i=1}^{N}\bm{y}_ilogP(\bm{p}_i|\bm{I}),
\end{equation}
where $\bm{y}_i$ is the one-hot representation for the true class of the $i^{th}$ instance, and $N$ represents the number of training instances.